\section{Research hypotheses}
\label{sec:hypotheses}

The following hypotheses are stated before the main experiments. Thresholds are
design choices for this project, not universal standards. A hypothesis is
falsified if the pre-specified comparison fails; null and negative results will
be reported.

\begin{table}[t]
\centering
\small
\caption{Pre-registered hypotheses and operational tests. All thresholds are
project-specific and may be revised only before results are inspected, with the
revision documented.}
\label{tab:hypotheses}
\begin{tabular}{@{}p{0.07\linewidth}p{0.52\linewidth}p{0.31\linewidth}@{}}
\toprule
ID & Hypothesis & Falsification criterion \\
\midrule
H1 & QLoRA adaptation improves held-out cybersecurity MCQ accuracy over the
unadapted base model by at least 5 percentage points at matched prompting. &
Observed improvement $<5$ points, or confidence interval includes zero. \\
H2 & At a fixed data and step budget, 4-bit QLoRA retains at least 90\% of the
held-out accuracy of bf16 LoRA while reducing peak training VRAM. &
Accuracy retention $<90\%$ or no measurable VRAM reduction. \\
H3 & Offline distillation from a larger local teacher improves held-out task
performance over supervised training on the same prompts at equal student
compute. & No improvement, or a degradation larger than the pre-specified
non-inferiority margin. \\
H4 & Retrieval augmentation reduces the lexical unsupported-token proxy on
evidence-grounded questions relative to closed-book generation, without a
significant drop in answer token F1. & No reduction, or token F1 drops beyond
the pre-specified margin. \\
H5 & FIM or prefix autocomplete produces usable completions (syntax-valid and
above a task-specific similarity threshold) at interactive latency on at least
one consumer profile. & Threshold not met on any target profile. \\
H6 & Narrow cybersecurity SFT causes a general-capability retention drop of no
more than 3 percentage points on the selected retention set. & Drop $>3$ points
with a confidence interval excluding the threshold. \\
\bottomrule
\end{tabular}
\end{table}

\paragraph{Why these hypotheses?}
H1--H2 address the core efficiency claim. H3 tests whether a larger local
teacher adds value beyond direct supervision. H4 tests whether retrieval
compensates for limited parametric knowledge. H5 addresses the autocomplete
objective, which is often omitted from cybersecurity model papers. H6 makes
catastrophic forgetting a reported outcome rather than a hidden cost
\citep{kirkpatrick2017}.

\paragraph{Analysis discipline.}
Each hypothesis has a primary metric, a fixed evaluation set version, a fixed
prompt template, and a pre-specified comparison. Secondary metrics are reported
but do not replace the primary test. Confidence intervals are computed by
bootstrap resampling \citep{efron1979bootstrap}. If a planned experiment cannot
be completed, the result is reported as missing rather than imputed.
